\chapter{Introdução}

A introdução de um relatório técnico de conclusão de curso constitui uma seção primordial que estabelece o contexto do projeto. Ela deve iniciar com uma exposição clara do tema central do trabalho, abordando a relevância do assunto dentro do campo da pesquisa proposta. Subsequentemente, é imprescindível delimitar o escopo do projeto, especificando os limites e o foco da pesquisa ou desenvolvimento realizado \cite{abnTeX2009}.
Os objetivos do projeto devem ser claramente delineados na introdução, fornecendo uma visão geral das metas que o trabalho almeja alcançar. Esses objetivos podem ser subdivididos em geral e específicos, se necessário, para uma compreensão mais detalhada das intenções do projeto. Detalhes sobre a metodologia empregada, os resultados e as conclusões não devem ser abordados nesta seção, pois serão tratados em partes específicas do relatório.
Ademais, a introdução pode incluir uma breve revisão da literatura relacionada ao tema, destacando trabalhos anteriores que tenham contribuído para a fundamentação teórica do projeto. Esse passo é importante pois auxilia na demonstração da base científica sobre a qual o trabalho, ainda que técnico, está construído e na justificativa da importância da pesquisa realizada.
Por fim, a introdução deve ser redigida de forma concisa e objetiva, visando capturar o interesse do leitor e fornecer uma visão geral clara do trabalho realizado. Ela serve como um prelúdio para as seções subsequentes que detalharão o projeto em profundidade.


\section{Objetivo Geral}

Indicar o objetivo geral.

\section{Objetivos Específicos}

Indicar cada objetivo específico do trabalho.

\section{Descrição da regra de negócio}

A descrição da regras de negocio é uma parte crucial do TCC, visto que elas definem como uma ferramenta ou produto de software deve operar para atender às necessidades dos usuários e às normas da organização. Uma descrição clara e detalhada garante que todos os stakeholders compreendam os requisitos e funcionalidades do software.


\section{Justificativa/Contextualização}

Esta subseção serve para explicar a importância e a relevância do projeto a ser realizado. Ela deve fornecer argumentos que demonstrem a necessidade do estudo, destacando o impacto potencial dos resultados na área de conhecimento ou na prática profissional. Essa subseção também contextualiza o problema abordado, situando-o em um quadro mais amplo, como tendências atuais, desafios específicos do setor ou lacunas na literatura existente. Ao justificar o trabalho, o autor estabelece a base para a compreensão do leitor sobre a importância do projeto e sua contribuição para a área de estudo ou aplicação prática.